\chapter{TEMPO FUOR DI SESTO}\label{tempo-fuor-di-sesto}

\hfill\break

Appresi la verità sullo Spin per la prima volta a una festa con gli
slittini, cinque anni dopo l'Evento di ottobre, in una freddissima e
pungente sera d'inverno. Fu Jason, al solito, che diede la notizia.

La serata cominciò con la cena dai Lawton. Jason era tornato a casa
dall'università per le vacanze di Natale, perciò era in un certo senso
un'occasione speciale, anche se ``solo in famiglia'' - ero stato
invitato dietro insistenza di Jase, probabilmente malgrado le obiezioni
di E.D.

«Sarebbe dovuta esserci anche tua madre» sussurrò Diane quando mi aprì
la porta. «Ho tentato di farla invitare da E.D. ma\ldots» Si strinse
nelle spalle.

«Non c'è problema,» la rassicurai, «Jason è già passato a salutarla.
Comunque, non si sente bene.» Era a letto col mal di testa, stranamente,
e io non ero proprio nella posizione di lamentarmi del comportamento di
E.D. Giusto il mese prima si era offerto di coprire le spese della
facoltà di medicina se avessi superato la prova di ammissione, «Perché
tuo padre ne sarebbe stato felice.» Era un gesto generoso e allo stesso
tempo sentimentalmente falso ma non potevo permettermi di rifiutare.

Marcus Dupree, mio padre, era stato l'amico più intimo (alcuni dicevano
l'unico) di E.D. Lawton ai tempi di Sacramento, quando cominciavano a
promuovere dispositivi di monitoraggio mediante aerostati al servizio
meteorologico e alla Polizia di frontiera. I miei ricordi di lui erano
vaghi e si erano fusi con le storie che raccontava mia madre, anche se
ricordavo distintamente i colpi alla porta la sera in cui morì. Era
l'unico figlio di una famiglia indigente franco-canadese del Maine,
orgoglioso del suo titolo di ingegnere, di talento ma ingenuo in fatto
di soldi: aveva perso i suoi risparmi in una serie di scommesse in
borsa, lasciando mia madre con un mutuo che non poteva sostenere.

Carol e E.D. assunsero mia madre come governante quando si trasferirono
a Est, forse un tentativo di E.D. di tenere in vita la memoria del suo
amico. Era importante che non avesse mai smesso di ricordare a mia madre
di averle fatto un favore? Che da allora in poi l'avesse trattata come
un accessorio di casa? Che conservasse una sorta di sistema delle caste
in cui la famiglia Dupree era considerata visibilmente di seconda
classe? Forse sì, forse no. «La generosità, di qualsiasi genere, è un
animale raro» diceva mia madre. Perciò forse ero io a immaginare (o ero
troppo sensibile all'argomento) il piacere che sembrava trarre dal
divario intellettivo tra me e Jason, la sua evidente convinzione che
fossi nato per fargli da comprimario, che fossi un normalissimo
parametro a confronto del quale poteva essere misurata la
straordinarietà di suo figlio.

Per fortuna sia io che Jase sapevamo che erano tutte stronzate.

Quando mi misi a sedere, Diane e Carol erano già a tavola. In via del
tutto eccezionale, quella sera Carol era sobria o almeno non tanto
ubriaca da notarsi. Un paio di anni prima aveva smesso di esercitare la
professione medica e in quei giorni tendeva a restare in casa per
evitare il rischio di essere accusata di guida in stato di ebbrezza. Mi
fece un sorriso di circostanza. «Tyler» disse. «Benvenuto.»

Dopo qualche minuto Jason e suo padre ci raggiunsero di sotto,
scambiandosi sguardi accigliati: era evidente che fosse successo
qualcosa. Quando prese posto accanto a me, Jase annuì distratto.

Come quasi tutte le occasioni di ritrovo della famiglia Lawton, la cena
fu cordiale ma tesa. Ci passammo i piselli e scambiammo due chiacchiere.
Carol era distante, E.D. stranamente silenzioso. Diane e Jason tentavano
di fare conversazione ma era chiaro che fosse successo qualcosa tra
Jason e suo padre di cui nessuno dei due voleva discutere. Jase sembrava
così controllato che arrivati al dolce mi chiesi se stesse male - non
aveva quasi staccato gli occhi dal piatto che a malapena aveva toccato.
Quando arrivò il momento di uscire per andare alla festa sugli slittini,
si alzò con evidente riluttanza e sembrò sul punto di rifiutare l'invito
finché E.D. Lawton lo rassicurò: «Fai pure, prenditi una serata libera.
Ti farà bene.» E io mi chiesi: ``Una serata libera da \emph{cosa}?''

Andammo alla festa con l'auto di Diane, una piccola Honda senza pretese,
una macchina del tipo ``la mia prima macchina'', come a Diane piaceva
descriverla. Mi sedetti dietro al posto del guidatore; Jase montò sul
sedile del passeggero accanto a sua sorella, ancora incupito, pigiando
con le ginocchia il cassetto del cruscotto.

«Cosa ti ha fatto,» chiese Diane, «ti ha sculacciato?»

«Macché!»

«Ti comporti come se lo avesse fatto.»

«Davvero? Mi dispiace.»

Il cielo naturalmente era buio. Mentre ci dirigevamo verso nord, i fari
scorrevano veloci sui prati innevati e su una muraglia di alberi senza
foglie. Tre giorni prima c'era stata una nevicata da record, seguita da
un'ondata di freddo che aveva imbalsamato la neve sotto una cappa di
ghiaccio, lì dove gli spazzaneve non erano passati. Qualche auto ci
superava con prudenza.

«Allora che è successo,» insistette Diane, «qualcosa di grave?»

Jason si strinse nelle spalle.

«Guerra? Peste? Carestia?»

Lui fece di nuovo spallucce e sollevò il collo della giacca.

\separatorefiori

Alla festa non andò molto meglio, per lui. Del resto, non è che fosse
proprio una festa.

Era un raduno di ex compagni di classe e conoscenti di Diane e Jason ai
tempi della Rice, ospitato dalla famiglia di un altro ex studente di
ritorno a casa da una qualche università della Ivy League. I suoi
genitori avevano tentato di organizzare un decoroso evento a tema:
tramezzini, cioccolata calda e discese in slitta sul lieve pendio dietro
casa. Ma per quasi tutti gli ospiti - tristi figli di papà che avevano
sciato a Zermatt o Gstaad molto prima di togliere l'apparecchio per i
denti - era solo un'altra scusa per bere di nascosto. Fuori, sotto fili
di luci colorate, le fiaschette d'argento circolavano liberamente; nel
seminterrato un tizio di nome Brent vendeva pasticche di ecstasy.

Jason trovò una sedia in un angolo e si mise a sedere guardando torvo
chiunque sembrasse cordiale. Diane mi presentò a una ragazza con gli
occhi enormi di nome Holly e poi mi mollò lì. Holly attaccò con un
monologo su tutti i film che aveva visto negli ultimi dodici mesi. Mi
fece girare per la stanza per quasi un'ora, fermandosi di tanto in tanto
per agguantare del sushi da un vassoio. Quando si scusò per andare in
bagno, me la svignai verso il punto dove Jason se ne stava imbronciato e
lo pregai di uscire con me.

«Non sono dell'umore adatto per andare in slitta.»

«Neanch'io. Fammi solo un favore, okay?»

E così indossammo stivali e giacche e ci trascinammo fuori. La serata
era fredda e senza vento. Una mezza dozzina di studenti della Rice ci
guardò di traverso, accalcati a fumare sotto il portico e avvolti in una
nuvola di fumo. Seguimmo un sentiero nella neve finché non fummo più o
meno soli in cima a una collina bassa, e ci soffermammo a guardare
qualcuno che senza troppo entusiasmo scivolava sulla slitta tra il
bagliore circense delle luci di Natale. Raccontai a Jason di Holly, che
mi si era appiccicata come una sanguisuga in abiti GAP. Si strinse nelle
spalle e disse: «Tutti abbiamo i nostri problemi.»

«Che cavolo hai stasera?»

Ma prima che potesse rispondere mi squillò il cellulare. Era Diane che
era tornata alla casa.

«Dove siete finiti? Holly è un po' incavolata. Abbandonarla in quel
modo\ldots{} sei stato proprio sgarbato, Tyler.»

«Sono sicuro che troverà qualcun altro su cui rovesciare i suoi
discorsi.»

«È solo nervosa. Non conosce quasi nessuno qui.»

«Scusa, ma perché dovrebbe riguardarmi?»

«Pensavo soltanto che voi due poteste legare.»

Battei le palpebre, sorpreso. «Legare?» Non è che si potesse
interpretarlo in maniera positiva. «Cosa cerchi di dirmi, che volevi
sistemarmi con lei?»

Rimase in silenzio per uno o due secondi accusatori. «Dai, Tyler\ldots{}
non prenderla così.»

Per cinque anni Diane era rimasta più o meno sfocata come un film
amatoriale, o così almeno mi era parso. C'erano stati dei momenti,
soprattutto dopo la partenza di Jason per l'università, in cui mi ero
sentito il suo migliore amico. Lei chiamava e parlavamo; andavamo in
giro a fare spese o vedevamo dei film assieme. Eravamo amici. Stretti.
Se c'era stata una qualche tensione sessuale tra di noi, era stata solo
da parte mia e io ero stato attento a non farla vedere, perché anche
quella parziale intimità era fragile - lo sapevo senza bisogno di
sentirmelo dire. Qualunque cosa Diane volesse da me, di certo non
includeva la passione.

Naturalmente E.D. non avrebbe mai tollerato una relazione tra me e Diane
a meno che non fosse controllata, essenzialmente infantile e priva di
ogni rischio di prendere una piega inaspettata. Ma la distanza tra noi
sembrava andar bene anche a Diane, e potevano passare anche mesi senza
quasi vederla. Riuscivo a salutarla da lontano quando aspettava
l'autobus della Rice (quando la frequentava ancora); ma in quei periodi
non mi chiamava e, nelle rare occasioni in cui raccoglievo il coraggio
sufficiente per chiamarla io, non era mai dell'umore giusto per parlare.

A quei tempi di tanto in tanto uscivo con delle ragazze della scuola, di
solito timide, che avrebbero preferito (spesso me lo dicevano
esplicitamente) frequentare un ragazzo più popolare ma si erano
rassegnate a una vita sociale di serie B. Nessuna di quelle relazioni
durò a lungo. A diciassette anni persi la verginità con una ragazza
carina e altissima, di nome Elaine Bowland; cercai di convincermi di
esserne innamorato ma dopo otto o nove settimane ci allontanammo con
rimpianto misto a sollievo.

Inaspettatamente, dopo ognuno di questi episodi, Diane mi chiamava e
parlavamo, io non facevo cenno a Elaine Bowland (o Toni Hickock o Sarah
Burstein) e Diane non aveva quasi mai il tempo di raccontarmi come aveva
trascorso le ore libere durante il nostro periodo di pausa, ma andava
bene così perché dopo poco tornavamo nella bolla, sospesi tra
romanticismo e finzione, infanzia e maturità.

Cercavo di non aspettarmi di più. Ma non riuscivo a smettere di
desiderare la sua compagnia. E credevo che lei desiderasse la mia.
Dopotutto, continuava a cercarla. Avevo notato com'era rilassata quando
ero con lei, il suo sorriso spontaneo quando entravo in una stanza,
quasi una dichiarazione: ``Ah, bene, Tyler è qui. Non succede mai niente
di brutto quando Tyler è qui.''

«Tyler?»

Mi chiedevo cosa potesse avere detto a Holly. ``Tyler è davvero
simpatico, ma sono anni che mi sta attaccato ai piedi come un
cane\ldots{} voi due sareste fantastici insieme!''

«Tyler?» Sembrava afflitta. «Tyler, se non vuoi parlare\ldots»

«In effetti non credo di volerlo.»

«Allora ti dispiace passarmi Jason?»

Gli diedi il cellulare. Jason ascoltò per un attimo, poi disse: «Siamo
sulla collina. No. No. Perché non vieni qui? Non fa così freddo. No.»

Non volevo vederla. Feci per andarmene. Jason mi lanciò il telefono e
disse: «Non fare lo stronzo, Tyler. Devo parlare sia a te che a Diane.»

«Di cosa?»

«Del futuro.»

Era un'osservazione fastidiosamente criptica. «Forse tu non hai freddo
ma io sì.» Stavo congelando.

«Questa cosa è più importante di qualunque problema tu abbia con mia
sorella.» Sembrava quasi comicamente serio. «E so cosa lei significa per
te.»

«Non significa niente per me.»

«Questo non sarebbe vero neanche se foste solo amici.»

«Noi \emph{siamo} solo amici.» Non gli avevo mai parlato davvero di
Diane; era uno degli argomenti da non sfiorare nelle nostre
conversazioni. «Chiediglielo tu stesso» aggiunsi.

«Sei incazzato perché ti ha presentato a questa Holly.»

«Non voglio parlarne.»

«Ma è solo perché ora gioca a fare la santa. Adesso si è fissata con
questa novità. Ha iniziato a leggere quei libri.»

«Quali libri?»

«Teologia dell'Apocalisse. In genere stanno nell'espositore dei best
seller. Del tipo: C.R. Ratel, \emph{Pregare al buio}, l'abnegazione
dell'io mondiale. Non voleva offenderti. È solo una specie di buona
azione.»

«Quindi questo giustifica tutto?» Mi allontanai ancora di qualche passo,
verso la casa. Cominciai a chiedermi come tornare a casa senza un
passaggio.

«Tyler» mormorò, e c'era qualcosa nella voce che mi fece voltare.
«Tyler. Ascolta. Mi hai chiesto cosa avessi.» Sospirò. «E.D. mi ha
raccontato delle cose dell'Evento di ottobre. Non sono ancora pubbliche.
Ho promesso che non ne avrei parlato. Ma voglio rompere quella promessa,
e lo faccio perché ci sono solo tre persone al mondo che per me
rappresentano la famiglia, una di queste è mio padre, e le altre due
siete tu e Diane. Perciò puoi sopportarmi ancora qualche altro minuto?»

Intravidi Diane, che a fatica saliva il pendio e cercava di infilarsi il
parka bianco da neve, con un braccio dentro e uno fuori.

Guardai la faccia di Jason, terribilmente triste alla luce confusa e
festiva proveniente dal basso. Quell'immagine mi spaventò e, nonostante
i miei sentimenti, acconsentii ad ascoltare ciò che aveva da dire.

\separatorefiori

Quando Diane raggiunse il gazebo, le sussurrò qualcosa. Lo guardò con
gli occhi sgranati e rimase a distanza da entrambi. Poi Jason cominciò a
parlare in maniera sommessa, metodica, quasi calmante, dispensando un
incubo come se fosse una ninna nanna.

Naturalmente aveva sentito tutto da E.D.

E.D. aveva fatto la sua fortuna dopo l'Evento di ottobre. Quando erano
venuti a mancare i satelliti, la Compagnia Lawton si era fatta avanti
con dei progetti per una tecnologia di rimpiazzo pratica e immediata:
aerostati ad alta quota, palloni sofisticati progettati per librarsi
senza meta nella stratosfera. Cinque anni dopo, gli aerostati di E.D.
trasportavano carichi e ripetitori per le telecomunicazioni, facendo
broadcast di dati e comunicazioni multipunto, in sostanza quasi tutto
quello che avrebbe potuto fare un satellite convenzionale (a parte il
GPS e l'astronomia). Il potere e l'influenza di E.D. erano cresciuti in
fretta. Più di recente aveva costituito un gruppo di pressione nel
settore aerospaziale, la Fondazione Perielio, e fornito consulenza al
Governo federale su alcuni progetti meno pubblici - in questo caso, il
programma VRA (Veicolo di Rientro Automatizzato) della NASA.

La NASA aveva affinato le sonde spaziali VRA da un paio d'anni. I lanci
iniziali erano stati messi a punto per fare indagini sullo scudo di
ottobre. Poteva essere penetrato? E dall'esterno potevano essere
recuperati dati utili?

Il primo tentativo fu quasi letteralmente un salto nel buio, un semplice
carico VRA in cima a un ammodernato Atlas IIAS della Lockheed Martin,
lanciato nell'oscurità assoluta sopra la base aerea di Vandenberg.
All'inizio sembrò un fallimento. Il satellite, progettato per
trascorrere una settimana in orbita, cadde nell'oceano Atlantico al
largo delle Bermuda qualche attimo dopo il lancio. «Come se,» disse
Jason, «avesse colpito il confine dell'Evento e fosse rimbalzato
indietro.» Ma \emph{non} era rimbalzato. «Quando lo hanno recuperato, ci
hanno trovato archiviati i dati di un'intera settimana.»

«Com'è possibile?»

«Non si tratta di cosa è \emph{possibile} ma di cosa è \emph{successo}.
Ed è \emph{successo} che il carico utile è stato sette giorni in orbita
ed è tornato la notte stessa in cui è partito. Sappiamo che è andata
proprio così perché è successo per ogni lancio, e di lanci ne hanno
fatti diversi.»

«\emph{Cosa} è successo? Di che cosa stai parlando, Jase? Di viaggi nel
tempo?»

«No\ldots{} non proprio.»

«Non \emph{proprio}?»

«Lascialo parlare» mi redarguì Diane con tono pacato.

«Ci sono indizi di ogni sorta su ciò che sta realmente accadendo» spiegò
Jason. «L'osservazione da terra sembra suggerire che i razzi, penetrando
la barriera, effettivamente accelerino prima di svanire, come se ne
fossero attratti. Eppure, i dati recuperati a bordo non mostrano quel
risultato. Le due serie di osservazioni non coincidono. Visti da qui, i
satelliti accelerano nella barriera e poi cadono quasi immediatamente;
gli stessi satelliti segnalano di essere avanzati senza difficoltà nelle
orbite programmate e di essere rimasti lì per il tempo assegnato, per
poi tornare sfruttando la loro stessa potenza settimane o mesi dopo.»
(``Come il cosmonauta russo,'' pensai, l'uomo la cui storia, mai
confermata o smentita ufficialmente, era diventata una specie di
leggenda metropolitana.)

Supponendo che entrambi gli insiemi di dati fossero legittimi, c'era
solo una spiegazione: il tempo oltre la barriera scorreva in modo
differente. Oppure, per invertire l'equazione, il tempo sulla Terra
scorreva più lentamente che nell'universo in generale.

«Capisci che significa?» domandò Jason. «All'inizio pensavamo di essere
in una specie di gabbia elettromagnetica che regola l'energia che
raggiunge la superficie della Terra. Ed è vero. Ma è solo un effetto
collaterale, una piccola parte di un disegno molto più grande.»

«Effetto collaterale di cosa?»

«Di ciò che chiamano un \emph{gradiente temporale}. Capisci l'importanza
che ha? Per ogni secondo che passa sulla Terra, oltre la barriera passa
un sacco di tempo in più.»

«Non ha senso» risposi immediatamente. «E con che cavolo di legge fisica
lo spiegheresti?»

«Gente con molta più esperienza di me sta cercando risposte a questa
domanda. Ma l'idea di un gradiente del tempo ha un certo potere
esplicativo. Se c'è un differenziale di tempo tra noi e l'universo, la
radiazione ambientale che raggiunge la superficie della Terra in un dato
momento - luce solare, raggi x, radiazione cosmica - verrebbe accelerata
in maniera proporzionale. E un anno di luce del Sole condensata in dieci
secondi sarebbe letale all'istante. Perciò la barriera elettromagnetica
intorno alla Terra non ci sta nascondendo, ci \emph{protegge}. Ci sta
riparando da tutta quella radiazione concentrata e, credo, spostata
verso il blu.»

«La finta luce solare» aggiunse Diane, che aveva afferrato il senso del
discorso.

«Esatto. Ci hanno dato la luce finta perché quella vera sarebbe stata
mortale. Giusto la quantità sufficiente e distribuita in modo
appropriato per imitare le stagioni, consentirci di coltivare la terra e
determinare le condizioni atmosferiche. Le maree, la traiettoria intorno
al Sole - massa, momento, gravitazione - \emph{tutto} questo viene
manipolato, non solo per rallentare il nostro moto ma anche per tenerci
in vita.»

«\emph{Manipolato}» ripetei. «Quindi non è un atto della natura. Si
tratta d'ingegneria.»

«Penso che dovremmo ammetterlo,» confermò Jason, «sì.»

«Lo stanno facendo \emph{per} noi.»

«Parlano di un'ipotetica intelligenza controllante.»

«Ma qual è lo scopo, cosa si dovrebbe ottenere?»

«Non lo so. Nessuno lo sa.»

Diane fissava il fratello attraverso uno spazio d'aria invernale fredda
e immobile. Si strinse nel parka e tremò. Non per la temperatura ma
perché era giunta alla domanda fondamentale: «\emph{Quanto} tempo,
Jason? \emph{Quanto} tempo passa là fuori?»

Là fuori oltre il vuoto del cielo.

Jason esitò, visibilmente restio a risponderle.

«Tanto tempo» ammise.

«Diccelo e basta» lo incalzò con tono flebile.

«Be', sono stati fatti tutti i tipi di misurazione. Ma nell'ultimo
lancio hanno fatto rimbalzare un segnale di calibrazione contro la
superficie della Luna. Ogni anno la Luna si allontana sempre di più
dalla Terra, lo sapevi? Di pochissimo ma quel poco è misurabile. Se si
misura quella distanza, si ottiene una specie di calendario rudimentale,
tanto più accurato quanto più tempo è trascorso. Aggiungi questo ad
altri segni, come il moto delle stelle vicine\ldots»

«\emph{Quanto tempo}, Jason?»

«Dall'Evento di ottobre sono passati cinque anni e un paio di mesi.
Oltre la barriera questo si traduce in poco più di cinquecento milioni
di anni.»

Era un numero che lasciava senza fiato.

Non riuscivo a pensare a nulla da dire. Neanche una parola. Mi aveva
lasciato incapace di esprimermi e di pensare. In quel momento non si
udiva neanche un suono, nient'altro che il vuoto frizzante della notte.

Poi Diane, che aveva già compreso l'essenza spaventosa del tutto,
bisbigliò: «E quanto tempo ci è rimasto?»

«Non so neanche questo. Dipende. In un certo senso siamo protetti dalla
barriera, ma quanto è efficace questa protezione? Ci sono comunque delle
realtà inevitabili. Il Sole è destinato a morire, come qualsiasi altra
stella. Brucia l'idrogeno, si espande e diventa più caldo man mano che
invecchia. La Terra si trova in una specie di zona abitabile del sistema
solare, e quella zona si sposta costantemente verso l'esterno. Ripeto,
siamo protetti, per il momento siamo a posto a prescindere da tutto. Ma
alla fine la Terra si troverà dentro l'eliosfera del Sole. Ne sarà
inghiottita. Passato un certo punto, non si può più tornare indietro.»

«Fra \emph{quanto}, Jase?»

La guardò compassionevole. «Quaranta, forse cinquant'anni, più o meno.»